\documentclass[UTF8]{ctexart}
\usepackage{xeCJK}
\usepackage{geometry}
\geometry{left=2cm,right=2cm,top=2.5cm,bottom=2.5cm}
\setlength{\headheight}{12.7pt}

\usepackage{titlesec}
\titleformat{\section}
  {\normalfont\large\bfseries}
  {\thesection}
  {1em}
  {}
\titlespacing*{\section}
  {0pt}
  {3.5ex plus 1ex minus .2ex}
  {2.3ex plus .2ex}

\usepackage{amsmath}
\usepackage{amssymb}
\usepackage{amsthm}
\usepackage{enumitem}
\setlist[enumerate]{itemsep=0pt,topsep=0pt,parsep=0pt,leftmargin=0pt,itemindent=2.5em}
\setlist[itemize]{itemsep=0pt,topsep=0pt,parsep=0pt,leftmargin=0pt,itemindent=2.5em}
\usepackage{fancyhdr}
\pagestyle{fancy}
\fancyhf{}
\usepackage{graphicx}
\usepackage{hyperref}
\usepackage{indentfirst}
\usepackage{lmodern}


\begin{document}
\setlength{\abovedisplayskip}{1pt}
\setlength{\belowdisplayskip}{1pt}

\section{并行赋值语句}

\hypertarget{sec:定义1.1}{}
\noindent\textbf{定义1.1 (并行赋值语句)}

对两两不同的变量$u_1,\cdots,u_n$ ($n\geqslant1$)和同类型的项
$t_1,\cdots,t_n$, 加入\textbf{并行赋值语句}：
\[
    u_1,\cdots,u_n:=t_1,\cdots,t_n.
\]

\vspace{10pt}

\hypertarget{sec:定义1.2}{}
\noindent\textbf{定义1.2 (并行赋值语句的语义)}

定义并行赋值语句的一步转移:
\[
    \langle u_1,\cdots,u_n:=t_1,\cdots,t_n,\sigma\rangle
    \overset{\mathrm{def}}{\to}
    \langle\mathrm{E},\sigma[u_1:=\sigma(t_1),\cdots,u_n:=\sigma(t_n)]\rangle,
\]
所有右端项在旧状态$\sigma$下同时求值更新. 显然
\[
    [\![u_1,\cdots,u_n:=t_1,\cdots,t_n]\!](\sigma)
    =\sigma[u_1:=\sigma(t_1),\cdots,u_n:=\sigma(t_n)].
\]

\vspace{10pt}

\hypertarget{sec:定义1.4}{}
\noindent\textbf{定义1.4 (并行赋值公理)}

把assignment公理推广为\textbf{并行赋值公理}(同时替换形式):
\[
    \{\,P[\vec{u}:=\vec{t}]\,\}\,u_1,\cdots,u_n:=t_1,\cdots,t_n\,\{P\},
\]
\[
    [\,P[\vec{u}:=\vec{t}]\,]\,u_1,\cdots,u_n:=t_1,\cdots,t_n\,[P],\\[-3pt]
\]
其中$P[\vec{u}:=\vec{t}]$是$P$的一次性同时替换.





\vspace{30pt}

\section{守卫语句}

\hypertarget{sec:定义2.1}{}
\noindent\textbf{定义2.1 (守卫语句)}

对$B:\mathrm{Bool}$和$S\in\mathcal{C}$, 加入\textbf{守卫语句}:
$\mathbf{if}\,\,B\to S\,\,\mathbf{fi}\in\mathcal{C}$.

\vspace{10pt}

\hypertarget{sec:定义2.2}{}
\noindent\textbf{定义2.2 (守卫语句的语义)}

用$\textbf{fail}$表示失败状态, 守卫语句的一步转移是:
\begin{enumerate}
    \item 当$\sigma\models B$时, $\langle\mathbf{if}\,\,B\to S
    \,\,\mathbf{fi}\rangle\to\langle S,\sigma\rangle$,
    \item 当$\sigma\not\models B$时, $\langle\mathbf{if}\,\,B\to S
    \,\,\mathbf{fi}\rangle\to\langle\mathrm{E},\textbf{fail}\rangle$.
\end{enumerate}

\vspace{10pt}

\hypertarget{sec:定义2.3}{}
\noindent\textbf{定义2.3 (守卫语句规则)}

守卫语句的部分正确性和完全正确性公式的推导规则是:
\vspace{5pt}
\[
    \mathrm{PW}:\,\frac{\{P\land B\}\,S\,\{Q\}}
    {\{P\}\,\mathbf{if}\,\,B\to S\,\,\mathbf{fi}\,\{Q\}},\quad
    \mathrm{TW}:\,\frac{P\to B,\,\{P\}\,S\,\{Q\}}
    {\{P\}\,\mathbf{if}\,\,B\to S\,\,\mathbf{fi}\,\{Q\}}.
\]

\end{document}
